\documentclass[12pt]{article}
\usepackage{graphicx} % Required for inserting images
\usepackage[margin=0.5in]{geometry}
\usepackage{amsmath, amssymb}
\usepackage{mathrsfs}


\usepackage{soul}


% Dr. David Brown Commands
\newcommand{\ds}{\displaystyle}
\newcommand{\R}{\mathbb{R}}
\newcommand{\N}{\mathbb{N}}
\newcommand{\Q}{\mathbb{Q}}
\newcommand{\C}{\mathbb{C}}


% Commands to get script r
\usepackage{calligra}
\DeclareMathAlphabet{\mathcalligra}{T1}{calligra}{m}{n}
\DeclareFontShape{T1}{calligra}{m}{n}{<->s*[2.2]callig15}{}
\newcommand{\scripty}[1]{\ensuremath{\mathcalligra{#1}}}

% Unit commands
\newcommand{\degree}{\text{°}}
\newcommand{\unitSpeed}{\frac{\text{m}}{\text{s}}}
\newcommand{\unitAcceleration}{\frac{\text{m}}{\text{s}^2}}

% Constant commands
\newcommand{\avogadro}{6.022\times10^{23}}

\newcommand{\boltzmannConstK}{1.381\times10^{-23}\frac{\text{J}}{\text{K}}}
\newcommand{\boltzmannConstInEV}{8.617\times10^{-5}\frac{\text{eV}}{\text{K}}}
\newcommand{\planckConst}{6.626\times10^{-34}\text{J}\cdot\text{s}}

\newcommand{\atmP}{1.01325\times10^5\,\text{Pa}}

% Known Value Commands
\newcommand{\electronMass}{9.109\times10^{-31}\,\text{kg}}
\newcommand{\protonMass}{1.673\times10^{-27}\,\text{kg}}

% Commands to easily write partial derivatives
\newcommand{\partDeriv}[2]{\frac{\partial #1}{\partial #2}}
\newcommand{\doublePartDeriv}[2]{\frac{\partial^2 #1}{\partial^2 #2}}


\title{Problem 7.44}
\author{Gabriel Decker}


\begin{document}


\maketitle
\begin{flushleft}


In this problem, we will consider the total number of photons $N$ in thermal equilibrium in a box with volume $V$ at temperature $T$.
We will compare it to the entropy, which the text says is equal to the entropy up to a constant.\\~\\

\textbf{Part a}\\
First, let's calculate $N$.
Recall the Planck distribution, which gives the number of photons at a certain energy level, based on the Bose-Einstein distribution.
\begin{equation}
    \overline{n}_\text{Pl}=\frac{1}{e^{\epsilon/kT}-1}
\end{equation}\\~\\

The total number of photons is simply twice the sum of the $\overline{n}_\text{Pl}$ for each energy state $\epsilon$.
Twice because of the two possible polarizations (see page 291 of the text).
So, this implies the following.
$$N=2\sum_\epsilon\overline{n}_\text{Pl}$$
$$\implies N=2\sum_\epsilon\frac{1}{e^{\epsilon/kT}-1}$$\\~\\

As shown in the text, recall that the energy levels are degenerate.
\begin{equation}
    \epsilon=\frac{hcn}{2L}
\end{equation}
with $n=\sqrt{n_x^2+n_y^2+n_z^2}$.
Plugging this into our $N$ equation, we get the following.
$$\implies N=2\sum_{n_x}\sum_{n_y}\sum_{n_z}\frac{1}{e^{hcn/2LkT}-1}$$\\~\\

Like in the book, we can use the thermodynamic limit to turn this into an integral.
We can express this integral in three-dimensional $n$-space in spherical coordinates, with the valid states taking place in $\phi\in[0,\pi/2]$, $\theta\in[0,\pi/2]$, and $n\in[0,\infty]$.
Recall that the infinitesimal volume in spherical coordinates is $dV=n^2\sin(\theta)\,dn\,d\phi \,d\theta$.
Therefore, we get the following integral.
$$\implies N=2\cdot\int_0^{\pi/2}\int_0^{\pi/2}\int_0^\infty\frac{1}{e^{hcn/2LkT}-1}n^2\sin(\theta)\,dn\,d\phi \,d\theta$$
$$\implies N=\pi\int_0^\infty\frac{n^2}{e^{hcn/2LkT}-1}\,dn$$\\~\\

The integrand can be simplified by using the substitution: $x=hcn/2LkT$.
This implies that $n=2LkTx/hc$ and $dn=(2LkT/hc)\,dx$.
Plugging this in, we get the following.
$$\implies N=\pi\int_0^\infty\frac{4(LkT)^2x^2}{(hc)^2}\frac{1}{e^{x}-1}\frac{2LkT}{hc}\,dx$$
$$\implies N=8\pi L^3\left(\frac{kT}{hc}\right)^3\int_0^\infty\frac{x^2}{e^{x}-1}\,dx$$
$$\implies N=8\pi V\left(\frac{kT}{hc}\right)^3\int_0^\infty\frac{x^2}{e^{x}-1}\,dx$$
This integral has to be evaluated numerically.
The integrand is approximately 2.404114.
$$\implies N=8\pi V\left(\frac{kT}{hc}\right)^3\cdot(2.404114)$$\\~\\~\\


\textbf{Part b}\\
We can compare this to the given entropy value by computing the entropy per particle $S/N$.
Recall the entropy result from the book.
\begin{equation}
    S(T)=\frac{32\pi^5}{45}V\left(\frac{kT}{hc}\right)^3k
\end{equation}
We can express $S/N$ in terms of a multiple of $k$.
$$\frac{S}{N}=\frac{(32\pi^5/45)V(kT/hc)^3k}{8\pi V(kT/hc)^3\cdot(2.404114)}$$
$$\implies\frac{S}{N}=\frac{32\pi^4/45}{8\cdot(2.404114)}k$$
$$\implies\frac{S}{N}=\frac{4\pi^4/45}{2.404114}k$$
$$\implies\frac{S}{N}=3.60157k$$\\~\\


\textbf{Part c}\\
Now, let's calculate the number of photons per cubic meter at various temperatures.
We can simply plug in $V=1\,\text{m}^3$ to do this.
$$N=8\pi V\left(\frac{kT}{hc}\right)^3\cdot(2.404114)$$
$$\implies N=8\pi (1\,\text{m}^3)\left(\frac{kT}{hc}\right)^3\cdot(2.404114)$$
We can plug in the cosmic background radiation temperature $T=2.73\,\text{K}$, room temperature $T=300\,\text{K}$, and the temperature of a kiln $T=1500\,\text{K}$.\\
\begin{table}[h]
\begin{center}
\begin{tabular}{ |c|c| }
    \hline
    Temp (K) & Number of Photons Per Square Meter\\
    2.73 & $4.1306\times10^{8 }$ \\
    300  & $5.4814\times10^{14}$ \\
    1500 & $6.8517\times10^{16}$ \\
    \hline
\end{tabular}
\end{center}
\end{table}



\end{flushleft}

\end{document}
